% GENERATED by `ggen sync` from `doc-pack` ontology over
% `docs/thesis/math_manufacturing/rdf/thesis.ttl` — chapter/section/
% paragraph structure and prose live as RDF individuals, not hand-authored
% .tex. Do not edit by hand; edit the .ttl instance data and re-sync.
\chapter{Introduction}

\section{Thesis statement}


A mathematics corpus of 219 statement instances, drawn from seven source
papers of an ongoing thesis series, was passed through an automated
pipeline that lowers each statement into Lean~4 syntax and submits it to
the Lean kernel for admission or rejection. The result is a disjoint
partition: 179 statements kernel-verified, 4 kernel-rejected after
bounded retry, 18 correctly withheld because they transitively depend on
an out-of-scope statement, and 17 cross-paper restatements collapsed by
shared label before scheduling. Every number in that sentence is derived
live in Chapter~\ref{ch:empirical} from artifacts on disk, not asserted
from memory.


This thesis makes one narrow technical claim and one broad, explicitly
labeled speculative claim, and it does not let the two blend.


\begin{description}
\item[The narrow claim (Part~I, Chapters~\ref{ch:background}--\ref{ch:empirical}):]
this pipeline is not a new theoretical apparatus. It is a second working
instance of a pattern already present in this codebase — the
\emph{manufacturing morphism} $\Act = \muop(\Adm)$ that the project's own
prior work already names and already implements once, for a different
domain (RDF ontology $\to$ PDDL planning text). Naming the Lean pipeline
``math manufacturing'' is an act of recognition, not invention: the same
admit/refuse/receipt structure recurs, and Chapter~\ref{ch:mathmfg} makes
the correspondence explicit, term for term.


\item[The broad claim (Part~II, Chapter~\ref{ch:civilizational}):] if
verified-mathematics production can be organized as a pipeline with a
defect taxonomy, a throughput number, and a receipt for every unit
produced — the way physical manufacturing is — then something changes
about the economics of formal verification, mathematical research, and
mathematics education at a civilizational scale. This claim is
speculative. It is not derived from the 219-statement result; it is
motivated by it. Chapter~\ref{ch:civilizational} marks every claim of
this kind with a dedicated \textsc{Speculation} environment so that a
reader — or an adversarial reviewer — can mechanically distinguish
``this is what we built and verified'' from ``this is what we think it
might mean if it scales.''
\end{description}

\section{Why the distinction matters}


A prior draft of this argument, written as prose in a conversational
exchange rather than as a checked document, contained a real accounting
error: a partition of 219 statements into verified/rejected/blocked/
deduplicated categories that summed to 218, not 219, because two
originally-verified pilot statements had been silently dropped from a
consolidated receipts file during a merge step. The error was caught,
traced to its exact cause, and fixed by re-verifying the two missing
statements against the real Lean kernel and re-inserting their receipts
— documented in full in Section~\ref{sec:receipts-gap}. That episode is
retained in this thesis, not smoothed over, because it is itself evidence
for the thesis statement: a manufacturing discipline is distinguished
from a narrative one precisely by the fact that its errors are
mechanically detectable and its claims are falsifiable against artifacts
on disk. A thesis arguing for that discipline should be held to it.

\section{Structure of this document}


Chapter~\ref{ch:background} establishes the pre-existing formal machinery
this thesis builds on: the Chatman Equation, admission algebra, receipts,
and the first working manufacturing morphism (\texttt{src/mfg.rs}).
Chapter~\ref{ch:mathmfg} states the second instantiation --- math
manufacturing --- as a formal correspondence to the first, closing with
a synthesis section (Section~\ref{sec:composite-chain}) naming every
stage of the full corpus-to-PDF chain in one place, for a reader who
wants the whole chain without reassembling it from three chapters.
Chapter~\ref{ch:empirical} reports the empirical run: exact counts,
the two infrastructure bugs found and fixed during the run, and the
receipts-consolidation gap and its fix. Chapter~\ref{ch:civilizational}
is the speculative civilizational-impact argument. Chapter~\ref{ch:related}
situates this work against prior autoformalization efforts.
Chapter~\ref{ch:limitations} states limitations without euphemism.
Chapter~\ref{ch:conclusion} concludes.

\chapter{Background: the Chatman Equation and manufacturing morphisms}
\label{ch:background}


This chapter is entirely evidence-bound: every construct below is quoted
or paraphrased from artifacts already present in this codebase prior to
this thesis, cited by exact file path.

\section{The Chatman Equation}


The project's foundational equation, stated in \texttt{ORIGINAL\_REQUEST.md}
and the admission-algebra paper (\texttt{docs/thesis/01\_admission\_algebra.tex}),
is
\[
  \Act = \muop(\Adm), \qquad
  \Obs \xrightarrow{\ \adm\ } \Adm \xrightarrow{\ \muop\ } \Act \xrightarrow{\ \text{receipt}\ } \Rec .
\]
Reading left to right: a raw observation space $\Obs$ is passed through
an \emph{admission map} $\adm$ into an admitted space $\Adm$ (only
observations that pass admission continue); the \emph{manufacturing
morphism} $\muop$ maps admitted observations to actuation $\Act$; and
every actuation is paired with a \emph{receipt} $\Rec$ — a computed,
content-addressed record of what happened, never asserted-in by fiat.
$\muop$ is named, in the project's own vocabulary, the ``manufacturing
morphism.'' This thesis's title is a direct restatement of that name
applied to a new domain.

\section{Admission algebra: refusal as data}
\label{sec:admission-algebra}


\texttt{01\_admission\_algebra.tex} establishes that admission cannot be
decided by meaning alone (a Rice's-theorem-style undecidability result,
\texttt{thm:rice} in the corpus) and instead formalizes admission as an
algebra over a refusal monoid: refusal is data, not a side effect
(\texttt{ax:refusal}), with a bottom element $\Rfsl$ characterizing full
admission (\texttt{prop:bottom}) and monotonicity under composition. The
practical corollary, realized in code, is that every admission stage
either produces a value or a named, typed refusal — never a panic, never
a silent default. This is the same discipline this project's own
\texttt{CLAUDE.md} states as an invariant: ``every error is a typed
\texttt{Refusal} variant.''

\section{Receipts: computed, not asserted}


\texttt{02\_receipt\_cryptography.tex} formalizes the receipt chain: a
BLAKE3 content hash folded against a genesis value, so that a receipt for
step $n$ is a function of the receipt for step $n-1$ and the actuation at
step $n$, never independently fabricated. This is the same mechanism
used throughout this thesis's own evidence: every count in
Chapter~\ref{ch:empirical} is derived by re-running a query against a
file on disk, not copied from a prior claim.

\section{The first manufacturing morphism: \texttt{src/mfg.rs}}
\label{sec:mfg-first}


A concrete implementation of $\muop$ already exists in this codebase,
independent of the work in this thesis. \texttt{src/mfg.rs} is a
literal manufacturing pipeline: it loads a Turtle RDF ontology describing
planning domain facts (types, predicates, actions, a problem instance),
extracts a STRIPS8-profile intermediate representation via SPARQL,
enforces PDDL8 structural bounds (arity, conjuncts, and parameter counts
each $\le 8$, mirroring \texttt{wasm4pm\_compat::pddl}'s
\texttt{PDDL8\_MAX\_*} constants) as a hard admission gate, and emits
PDDL domain and problem text as its actuation. Its module doc comment
states the shape directly: ``RDF ontology $\to$ PDDL8 domain/problem
text.''


Its error type,
\begin{quote}
\texttt{enum MfgError \{ Graph(String), BoundExceeded \{ what, limit,
got, detail \}, Shape(String), Pddl8(String) \}}
\end{quote}
is a defect taxonomy for this pipeline: a graph-parse failure, a
structural-bound violation (with the exact limit and the exact value
that exceeded it), a missing or malformed ontology fact, or a
round-trip failure back through the downstream planner. Every one of
these is a typed refusal in the admission-algebra sense of
Section~\ref{sec:admission-algebra} above — none of them is a panic or a silent fallback.


\begin{definition}[$\muop$-pipeline]
\label{def:mu-pipeline}
A \emph{$\muop$-pipeline} is a concrete instantiation of
$\Act = \muop(\Adm)$ consisting of: (i) a raw input domain $\Obs$;
(ii) an admission map $\adm : \Obs \to \Adm \cup \{\Rfsl\} \times \Rec_f$
that either admits an input or produces a typed refusal; (iii) a
manufacturing morphism $\muop : \Adm \to \Act$ producing a concrete
actuation artifact; and (iv) a receipt function computing $\Rec$ from
the actuation, never asserted.
\end{definition}


\texttt{src/mfg.rs} instantiates Definition~\ref{def:mu-pipeline} with
$\Obs$ = RDF ontology text, $\adm$ = SPARQL extraction plus PDDL8 bound
enforcement, $\muop$ = domain/problem text emission, and $\Act$ = PDDL8
text validated by round-trip through \texttt{bcinr-pddl}. Call this
\textbf{$\muop$-pipeline~1}. Chapter~\ref{ch:mathmfg} instantiates the
same definition a second time, with mathematics statements as the input
domain and the Lean~4 kernel as the admission map — \textbf{$\muop$-pipeline~2}
— and states the correspondence between the two, term for term.

\chapter{Math manufacturing: the second manufacturing morphism}
\label{ch:mathmfg}

\section{The correspondence, term for term}


Table~\ref{tab:correspondence} states the correspondence between
$\muop$-pipeline~1 (\texttt{src/mfg.rs}, Section~\ref{sec:mfg-first}) and
$\muop$-pipeline~2, the autoformalization pipeline built and run this
thesis reports on.


\begin{table}[h]
\centering
\small
\begin{tabular}{@{}p{2.1cm}p{5.1cm}p{5.1cm}@{}}
\toprule
\textbf{Term} & \textbf{$\muop$-pipeline 1 (RDF $\to$ PDDL)} & \textbf{$\muop$-pipeline 2 (Math $\to$ Lean)} \\
\midrule
$\Obs$ (raw input) & Turtle ontology text & LaTeX statement text, corpus RDF \\
$\adm$ (admission map) & SPARQL extraction + PDDL8 bounds & Lean 4 kernel typecheck \\
$\Adm$ (admitted domain) & STRIPS8 intermediate representation & Lean 4 syntax draft (pre-kernel) \\
$\muop$ (manufacturing morphism) & domain/problem text emission & \texttt{.lean} source file emission \\
$\Act$ (actuation) & PDDL8 text, round-tripped & kernel-accepted \texttt{.lean} artifact \\
Defect taxonomy & \texttt{MfgError} (Graph/BoundExceeded/Shape/Pddl8) & verified / unformalized / blocked \\
$\Rec$ (receipt) & (implicit: round-trip check) & \texttt{formalization\_receipts.jsonl} line \\
\bottomrule
\end{tabular}
\caption{Term-for-term correspondence between the two working
$\muop$-pipeline instantiations in this codebase.}
\label{tab:correspondence}
\end{table}


The correspondence is not decorative. Both pipelines share the structural
property that admission is a hard, typed gate rather than a best-effort
check: $\muop$-pipeline~1 will not emit PDDL8 text that violates its
arity/conjunct/parameter bounds, and $\muop$-pipeline~2 will not record a
statement as \texttt{verified} without an actual, successful invocation
of the real \texttt{lean} binary against it (Section~\ref{sec:no-self-report}).

\section{A third instantiation: this document}
\label{sec:doc-mfg-correspondence}


Definition~\ref{def:mu-pipeline} generalizes over any ($\Obs$, admission
map, $\muop$, $\Act$, defect taxonomy, receipt) tuple; Section~\ref{sec:mfg-first}
and Chapter~\ref{ch:mathmfg} instantiate it twice. This document's own
Chapters~1--8 are a third instantiation, not a new mechanism: $\Obs$ is
the chapter/section/paragraph RDF graph in
\texttt{docs/thesis/math\_manufacturing/rdf/thesis.ttl}; the admission
map is \texttt{ggen sync}'s template rendering followed by a successful
\texttt{pdflatex} compile (a malformed paragraph or a broken cross-
reference fails one of these two gates rather than silently producing a
bad PDF); $\muop$ is the \texttt{doc-pack} template render itself; $\Act$
is \texttt{thesis.pdf}; the defect taxonomy is (missing label, register
mismatch, template render failure, LaTeX compile failure); and the
receipt is the \texttt{graph\_hash\_hex} value \texttt{ggen sync}'s own
JSON output already reports for every run.


That receipt is deliberately not quoted here as a literal value: doing
so would require this sentence to already contain the hash of a document
that includes this sentence, a fixed-point dependency this thesis does
not attempt to resolve. The receipt is instead operational rather than
quoted -- re-running \texttt{docs/thesis/math\_manufacturing/regenerate.sh}
after any edit to \texttt{thesis.ttl} deterministically reproduces a
\texttt{graph\_hash\_hex} in its JSON output, and that value changes if
and only if the RDF content changes, which is the property a receipt in
the sense of Section~\ref{sec:admission-algebra} is required to have.


\begin{finding}
This thesis document is itself an instance of the manufacturing-morphism
pattern it describes: its chapters are admitted RDF facts, not asserted
prose, rendered through the same admit/render/receipt shape as
Chapter~\ref{ch:mathmfg}'s Lean pipeline and Section~\ref{sec:mfg-first}'s
PDDL pipeline. This is a structural observation checkable by re-running
the regenerate script, not a Lean-checked or otherwise mechanically
proved claim -- it is reported at the same evidentiary standing as
Chapter~\ref{ch:empirical}'s other findings, no stronger.
\end{finding}

\section{The pipeline in detail}


\begin{definition}[Math manufacturing pipeline]
Given a corpus of statement individuals $s \in \Obs$, each carrying a
LaTeX statement string, a \texttt{kind} (axiom, definition, construction,
theorem, proposition, corollary, or lemma), and a set of
\texttt{dependsOn} edges to other statements:
\begin{enumerate}[nosep]
  \item \textbf{Scoping.} Statements sharing a label with an
  already-processed statement are deduplicated to the first occurrence
  (Section~\ref{sec:dedup}); a designated out-of-scope statement and its
  transitive dependents are excluded a priori (Section~\ref{sec:excluded}).
  \item \textbf{Scheduling.} The remaining statements are ordered into
  topological waves over the \texttt{dependsOn} DAG (Kahn's algorithm):
  wave~0 contains statements with no unresolved dependencies; wave~$n$
  contains statements whose dependencies were all resolved (verified or
  blocked) by the end of wave $n{-}1$.
  \item \textbf{Admission.} Within a wave, each statement is drafted into
  Lean~4 syntax by an agent given the statement's LaTeX text and the
  already-verified Lean declarations of its dependencies, then
  typechecked by the real \texttt{lean} kernel. On rejection, the exact
  kernel error is fed back to the same agent for revision, bounded to
  three attempts total.
  \item \textbf{Disposition.} A statement is recorded \texttt{verified}
  on kernel acceptance, \texttt{unformalized} after three failed
  attempts (with the last real kernel error retained), or
  \texttt{blocked} if any dependency did not resolve to
  \texttt{verified}.
  \item \textbf{Receipt.} Every disposition is appended as one line to
  \texttt{formalization\_receipts.jsonl}, regardless of outcome.
\end{enumerate}
\end{definition}


\subsection{No self-reported success}
\label{sec:no-self-report}


A structural requirement of this pipeline, inherited from the admission-
algebra discipline of Section~\ref{sec:admission-algebra}, is that an agent's own claim of
success is never trusted as the disposition. Every \texttt{verified}
disposition in Chapter~\ref{ch:empirical} is independently re-checked by
invoking the real \texttt{lean} kernel directly against the \texttt{.lean}
file on disk, outside of and after the agent's own report. This check
caught one concrete discrepancy during the run reported in
Chapter~\ref{ch:empirical}: a statement self-reported as
\texttt{unformalized} was found, on independent re-verification, to
type-check cleanly, and its disposition was corrected before being
recorded in the receipts file.


\subsection{Deduplication}
\label{sec:dedup}


The source corpus contains 219 statement individuals but only 202
distinct labels (Section~\ref{sec:corpus-numbers}): 17 labels are
restated across more than one source paper (the same concept,
independently re-derived in each paper's own argument). Because
\texttt{dependsOn} edges reference labels rather than per-paper
individuals, the scheduler resolves each label once, at its first
occurrence, rather than re-formalizing 17 statements that would produce
identical Lean declarations.


\subsection{A priori exclusion}
\label{sec:excluded}


One statement, \texttt{thm:faithful} (a faithfulness/projection theorem
with cryptographic and probabilistic content), is excluded from scope
before scheduling begins, together with every statement that transitively
depends on it. This is not a failure of the pipeline; it is a scoping
decision reported honestly, discussed further in
Chapter~\ref{ch:limitations}: bare Lean~4 core, with no mathlib
installed, is not the appropriate kernel for cryptographic/probabilistic
content, and formalizing it would require a different toolchain
(Chapter~\ref{ch:limitations} returns to this point).

\section{The composite chain: corpus, kernel, and document}
\label{sec:composite-chain}


Chapters~\ref{ch:background}--\ref{ch:empirical} and
Section~\ref{sec:doc-mfg-correspondence} each establish one stage of a
single chain, but no single place names all five stages together.
Table~\ref{tab:composite-chain} does that -- every fact in it is
restated, not new, and each row cites the section that establishes it.


Table~\ref{tab:composite-chain} names the stages; the chain can also
be written as a single sequence of admissions, in the same
$\Obs \to \Adm \to \Act$ notation Chapter~\ref{ch:background} already
defines:
\[
\begin{aligned}
O_{\mathrm{raw}}
&\xrightarrow{\ \alpha_{\mathrm{praxis}}\ } O^{*}_{\mathrm{praxis}}
\xrightarrow{\ \muop_{\mathrm{L4}}\ } A_{\mathrm{L4}}
\xrightarrow{\ \rho_{\mathrm{L4}}\ } R_{\mathrm{L4}} \\
&\xrightarrow{\ \muop_{\mathrm{ggen}}\ } A_{\mathrm{tex}}
\xrightarrow{\ \muop_{\mathrm{latex}}\ } A_{\mathrm{pdf}}
\xrightarrow{\ \rho_{\mathrm{doc}}\ } R_{\mathrm{doc}}
\end{aligned}
\]
where $O_{\mathrm{raw}}$ is the seven source papers' prose;
$\alpha_{\mathrm{praxis}}$ is corpus extraction and SHACL/citation/
notation validation; $O^{*}_{\mathrm{praxis}}$ is \texttt{corpus.ttl};
$\muop_{\mathrm{L4}}$ is the LLM-draft-plus-kernel-typecheck step;
$A_{\mathrm{L4}}$ is a verified \texttt{.lean} file; $\rho_{\mathrm{L4}}$
is the receipt function producing $R_{\mathrm{L4}}$ (one
\texttt{formalization\_receipts.jsonl} line); $\muop_{\mathrm{ggen}}$ is
\texttt{ggen sync}; $A_{\mathrm{tex}}$ is the generated
\texttt{chapters\_body.tex}; $\muop_{\mathrm{latex}}$ is the
\texttt{pdflatex} compiler; $A_{\mathrm{pdf}}$ is \texttt{thesis.pdf};
and $\rho_{\mathrm{doc}}$ produces $R_{\mathrm{doc}}$, the
\texttt{graph\_hash\_hex} value discussed in
Section~\ref{sec:doc-mfg-correspondence}. This is notation for the same
two Definition~\ref{def:mu-pipeline} instances already named in
Table~\ref{tab:composite-chain}, not a claim beyond them.


\begin{table}[h]
\centering
\small
\begin{tabular}{@{}p{2.1cm}p{2.6cm}p{2.6cm}p{2.3cm}p{2.3cm}p{2.3cm}@{}}
\toprule
\textbf{Stage} & \textbf{Input} & \textbf{Admission gate} & \textbf{Morphism} & \textbf{Artifact} & \textbf{Refusal mode} \\
\midrule
Praxis corpus & raw thesis prose, 7 papers & RDF parse + SHACL/citation/notation checks (Sec.~\ref{sec:corpus-numbers}) & \texttt{build\_corpus.py} extraction & \texttt{corpus.ttl} (219 IRIs, 202 labels) & malformed fact, dangling citation \\
Lean 4 & statement label + LaTeX text & kernel typecheck, no self-report (Sec.~\ref{sec:no-self-report}) & LLM draft + \texttt{lean} invocation & \texttt{.lean} file & kernel rejection after 3 attempts \\
Scheduler & \texttt{dependsOn} DAG & topological wave readiness & Kahn's-algorithm wave execution & wave disposition & blocked dependency \\
ggen & document RDF graph & template render + \texttt{pdflatex} compile (Sec.~\ref{sec:doc-mfg-correspondence}) & \texttt{ggen sync} (SPARQL + Tera) & \texttt{chapters\_body.tex} & missing label, register mismatch \\
LaTeX & generated \texttt{.tex} & \texttt{pdflatex} exit code & TeX compilation & \texttt{thesis.pdf} & compile error \\
\bottomrule
\end{tabular}
\caption{The five stages of the corpus-to-PDF chain. Every cell restates a fact established elsewhere in this thesis; none is introduced here.}
\label{tab:composite-chain}
\end{table}


A second table states what each layer does and, as importantly, what it
does not do -- a lawful division worth stating plainly, since conflating
any two of these is the most common way to overclaim this kind of
pipeline: it is neither true that ``the LLM proves the theorem'' nor
that ``ggen proves the math'' nor that ``LaTeX validates anything.''


\begin{table}[h]
\centering
\small
\begin{tabular}{@{}p{2.1cm}p{5.3cm}p{5.3cm}@{}}
\toprule
\textbf{Layer} & \textbf{Does} & \textbf{Does not do} \\
\midrule
Praxis corpus & holds the admitted semantic/document graph & does not itself kernel-verify any statement \\
Lean 4 & admits or refuses statements by kernel typecheck & does not judge whether informal prose was translated faithfully \\
ggen & manufactures artifacts from graph facts via SPARQL/Tera & does not invent or validate the graph's meaning \\
LaTeX & typesets the rendered artifact & does not validate semantic or mathematical truth \\
\bottomrule
\end{tabular}
\caption{Lawful division of labor across the chain: each layer's guarantee, and its explicit non-guarantee.}
\end{table}


This composite chain is not a fourth, novel mechanism requiring its own
definition. It is Definition~\ref{def:mu-pipeline} applied twice from a
shared source: the admitted Praxis corpus is the $\Obs$ that both the
Lean pipeline (Chapter~\ref{ch:mathmfg}) and the document pipeline
(Section~\ref{sec:doc-mfg-correspondence}) separately consume as their
own admission map's input. A single composite predicate asserting this
whole chain as one checked logical equivalence would be new, unchecked
formal machinery -- no SAT/SMT/Lean encoding of such a predicate exists
in this repository -- and this thesis declines to state one, for the
same reason Section~\ref{sec:doc-mfg-correspondence} already declines to
quote a literal receipt hash that would require a fixed point: an
unearned formal claim is a worse defect than an informal but accurate
one.


Every stage in Table~\ref{tab:composite-chain} is independently
replayable by a reader, not merely described:
\begin{itemize}[nosep]
\item \textbf{corpus query} -- re-run the SPARQL \texttt{SELECT} shown
in Section~\ref{sec:corpus-numbers} against \texttt{docs/thesis/rdf/corpus.ttl}.
\item \textbf{Lean recheck} -- \texttt{PATH="\$HOME/.elan/bin:\$PATH" lean
tools/paper-factory/lean-pilot/<statement>.lean} for any \texttt{verified}
statement (Section~\ref{sec:lean-admission}).
\item \textbf{ggen regenerate} -- \texttt{docs/thesis/math\_manufacturing/regenerate.sh}
reproduces \texttt{generated/chapters\_body.tex} and reports
\texttt{graph\_hash\_hex}.
\item \textbf{pdflatex compile} -- two passes of \texttt{pdflatex thesis.tex}
inside \texttt{docs/thesis/math\_manufacturing/} reproduce \texttt{thesis.pdf}.
\item \textbf{register-boundary check} -- the SPARQL query in
Section~\ref{sec:speculation-boundary} re-confirms the evidence-bound/
speculative partition below.
\end{itemize}
None of these five steps depends on trusting a prior run's report; each
can be executed independently against artifacts already in this
repository.

\section{Semantic Rice quarantine}
\label{sec:rice-quarantine}


Chapter~\ref{ch:background} recalls that admission-algebra's own
Rice's-theorem-style result (\texttt{thm:rice}) says admission cannot
decide meaning in general. This section states the corresponding
quarantine rule for this thesis's Lean pipeline precisely, so that
\texttt{LeanAdmitted} (Section~\ref{sec:lean-admission}) is never
mistaken for a stronger claim than it is.


\begin{definition}[Semantic Rice quarantine]
A semantic claim receives Lean~4 standing only when its meaning is
carried by an admitted formal object: a definition, an interpretation
map, an equivalence or refinement relation, or a preservation theorem,
each itself a statement in the corpus. Lean~4 proves consequences of
admitted definitions. It does not prove that an informal semantic
intention has been fully captured by those definitions, unless that
capture is itself formalized as one of the objects above.
\end{definition}
Under this rule, \texttt{thm:faithful} is excluded
(Section~\ref{sec:corpus-numbers}) precisely because its cryptographic
and probabilistic semantics are not admitted into the bare-Lean-core
surface this run uses: no interpretation map for those semantics exists
in the corpus, so no Lean file could stand for that theorem's intended
meaning rather than an unrelated, easier substitute.


\begin{table}[h]
\centering
\small
\begin{tabular}{@{}p{6.2cm}p{4cm}@{}}
\toprule
\textbf{Claim type} & \textbf{Status} \\
\midrule
Typed theorem over admitted definitions & Lean-admissible \\
Duplicate label (same concept, restated across papers) & quotient/dedup \\
Dependency on an excluded theorem & blocked \\
Cryptographic/probabilistic semantic claim without library support & Rice-quarantined \\
Informal intended meaning without an interpretation map & not Lean-admitted \\
\bottomrule
\end{tabular}
\caption{Classification of claim types under the semantic Rice quarantine.}
\label{tab:rice-quarantine}
\end{table}
No reader should equate ``Lean proves this formal file'' with ``Lean
proves every possible intended semantic interpretation of the original
prose.'' The former is what Table~\ref{tab:disposition} reports; the latter is not
claimed anywhere in this thesis.

\section{Hardening the admission gate: the praxis-lean crate}
\label{sec:praxis-lean}


After the empirical run reported in Chapter~\ref{ch:empirical}, the
admission principles of Section~\ref{sec:no-self-report} were hardened
from one-off shell invocations into a compiled, workspace-integrated
Rust crate, \texttt{praxis-lean} (binary \texttt{praxis-l4}), added to
this repository's Cargo workspace. Its core law, stated in its own
\texttt{lib.rs} doc comment, restates
Section~\ref{sec:no-self-report} as a checked type-level invariant
rather than a convention to remember:
\[
\mathrm{Verified}(s) \iff \mathrm{KernelAccepts}(s) \land \mathrm{NoSorry}(s) \land \mathrm{NoUnauthorizedAxiom}(s).
\]
Every admission or refusal outcome is a typed \texttt{LeanRefusal}
variant (\texttt{KernelRejected}, \texttt{SorryFound},
\texttt{UnauthorizedAxiom}, \texttt{OrphanLabel}, \texttt{OrphanFile},
\texttt{OrphanReceipt}, among others) -- never a panic, matching the
admission-algebra discipline of Section~\ref{sec:admission-algebra}.


The crate's CLI itself is manufactured, not hand-written: six
\texttt{praxis:CliCommand} instances (\texttt{init}, \texttt{verify},
\texttt{no-sorry}, \texttt{index-build}, \texttt{reconcile},
\texttt{report}, noun \texttt{l4}) were added to
\texttt{schema/praxis.ttl} and rendered into
\texttt{crates/praxis-lean/src/verbs/} by \texttt{ggen sync}, the same
mechanism Section~\ref{sec:doc-mfg-correspondence} already documents for
this thesis's own chapters -- a third artifact class (compiled CLI
routing, not LaTeX or Lean) manufactured from the same admitted RDF
graph.


\begin{finding}
Running \texttt{praxis-l4 l4 index-build} against the real
\texttt{corpus.ttl} independently reproduces 202 records -- exactly the
distinct-label count Section~\ref{sec:corpus-numbers} reports from a
wholly separate Python/SPARQL extraction path. Running
\texttt{praxis-l4 l4 reconcile} against that index and the real
\texttt{formalization\_receipts.jsonl} reports 202 labels, 201 receipt
lines, a single \texttt{missing\_receipts} entry
(\texttt{thm:faithful}, the one a-priori exclusion), a disposition
count of 179 verified plus 4 unformalized plus 18 blocked (summing to
201), and zero orphan receipts -- an exact, independent reconfirmation
of Table~\ref{tab:disposition}'s accounting invariant, produced by a compiled Rust tool
with no code path shared with the original extraction pipeline.
\end{finding}


Running \texttt{praxis-l4 l4 no-sorry} against all 183 real
\texttt{.lean} files under \texttt{tools/paper-factory/lean-pilot/}
independently reconfirms zero uses of \texttt{sorry} across the entire
verified set, and reports 245 real \texttt{axiom} declarations --
expected, since bare Lean~4 core (Section~\ref{ch:limitations}) has no
library to build primitives like \texttt{Obs} or \texttt{Halts} from
except axioms. Building this checker surfaced one real defect before
this reconfirmation was trustworthy: an earlier default policy also
forbade the bare token \texttt{admit}, which is not a Lean~4 tactic at
all (Lean~4's only real proof hole is \texttt{sorry}; \texttt{admit}
is a Coq idiom) -- it flagged five matches that were all the corpus's
own \texttt{LifeHom} type's constructor literally named \texttt{admit},
a real identifier, not a proof escape hatch. The policy default was
corrected before reporting the zero-\texttt{sorry} finding above, and
the corrected behavior is covered by a regression test.


The crate also proves out a second, additive Lean verification
lane: \texttt{praxis-l4 l4 init} scaffolds a real Lake package (a
\texttt{lean-toolchain}, \texttt{lakefile.lean}, and a
\texttt{lean\_lib} target) in a new sibling directory,
\texttt{tools/paper-factory/lean-lake/}, deliberately not touching any
of the 183 bare-\texttt{lean}-verified files. \texttt{lake build} and
\texttt{lake env lean Praxis.lean} both succeed against the real Lean
4.31.0/Lake 5.0.0 toolchain. Closing the gap that \texttt{reconcile} and
\texttt{report} currently accept the pre-existing receipts schema only
through a schema-tolerant compatibility reader (rather than the richer
schema~v2 record this crate defines, with per-file hashes and a
genesis-folded chain hash) remains explicit future work, not silently
assumed done.


This lane was then linked to Mathlib, and its prebuilt cache used
rather than a from-source build:
\texttt{require mathlib from git ... @ "v4.31.0"} (pinned to the one
Mathlib git tag whose own \texttt{lean-toolchain} exactly matches the
toolchain already installed, avoiding a second toolchain download), then
\texttt{lake exe cache get}, which downloaded and decompressed all
8{,}542 of Mathlib's prebuilt \texttt{.olean} files rather than
compiling roughly a thousand Mathlib modules from source. A trivial
\texttt{import Mathlib} file typechecked in 38 seconds wall-clock against
that cache -- a from-source build of Mathlib takes hours, so this is
direct evidence the cache, not a rebuild, is what made the lane usable.


\begin{finding}
Within this Mathlib-linked lane, four files were reformalized to compose
pre-built content instead of declaring axioms from scratch. \texttt{Bits256}
became the real \texttt{BitVec 256} rather than an opaque axiomatized
type; four of \texttt{def:receipt}'s six axiomatized fields
(\texttt{DenialWord}, \texttt{TransitionId}, \texttt{Version},
\texttt{RefusalReason}, \texttt{Fitness}) were composed from
\texttt{Nat}, \texttt{String}, \texttt{Prod}, and -- reusing this same
pilot's own \texttt{Deny} type from \texttt{prop:monoid} rather than a
second, unrelated axiom -- a dependent pair over it; \texttt{prop:monoid}'s
five hand-proved lattice lemmas were replaced by five one-line citations
of Mathlib's pre-built \texttt{BooleanAlgebra} instance on
\texttt{Fin n} $\to$ \texttt{Bool}; and \texttt{thm:rice}'s bare-core
version, which axiomatizes 13 facts about an abstract computability layer
and then hand-proves the halting-problem reduction (confirmed by direct
grep count), was replaced by a direct corollary of Mathlib's own
\texttt{ComputablePred.rice\_2} (Lean's actual identifier uses a Unicode
subscript~2, rendered here in plain ASCII) and
\texttt{ComputablePred.halting\_problem}
(\texttt{Mathlib.Computability.Halting}, formalizing Carneiro 2019) --
citing an existing, published, from-first-principles Lean proof of
Rice's theorem rather than re-deriving it. Lean's own
\texttt{\#print axioms} confirms the resulting corollaries depend only
on Lean's three standard foundational axioms
(\texttt{propext}, \texttt{Classical.choice}, \texttt{Quot.sound}) --
zero corpus-specific custom axioms, down from 13 in the bare-core file.
Independently re-running \texttt{praxis-l4 l4 no-sorry} against this
lane (Section~\ref{sec:praxis-lean}) confirms the count directly: 5
axioms remain across all four files (\texttt{chainH}, \texttt{chainStep}
-- the genuinely irreducible cryptographic primitives -- plus
\texttt{Obs}, \texttt{Sim}, and one bundled \texttt{Sim\_equiv} axiom,
down from 3 separate hand-written equivalence axioms).
\end{finding}


This pilot covers 4 of the corpus's 219 statements. It is evidence
that the Mathlib lane works and that real reuse is possible, not a claim
that the remaining 215 statements have been reformalized -- that
migration, and closing the receipts-schema compatibility gap noted
above, remain future work, exactly as Section~\ref{ch:limitations}
states.

\chapter{Empirical results}
\label{ch:empirical}


Every number in this chapter was re-derived at drafting time, not copied
from earlier conversational prose, by two commands: a Python pass
counting the \texttt{status} field across every line of
\texttt{tools/paper-factory/lean-pilot/formalization\_receipts.jsonl},
and a SPARQL query against \texttt{docs/thesis/rdf/corpus.ttl} selecting
every \texttt{math:Statement} individual's label and kind. Both are
reproducible by any reader with the corpus and receipts file in hand.

\section{Corpus numbers}
\label{sec:corpus-numbers}


\begin{table}[h]
\centering
\begin{tabular}{@{}lr@{}}
\toprule
\textbf{Quantity} & \textbf{Count} \\
\midrule
Statement individuals (IRIs) in corpus & 219 \\
Distinct statement labels & 202 \\
Duplicate-label IRIs (restated across papers, deduplicated) & 17 \\
Proof individuals & 95 \\
Distinct symbols (post notation-conflict merge) & 53 \\
\bottomrule
\end{tabular}
\caption{Corpus-level counts, re-derived by SPARQL against
\texttt{docs/thesis/rdf/corpus.ttl}.}
\end{table}

\section{Formalization disposition}


\begin{table}[h]
\centering
\begin{tabular}{@{}lr@{}}
\toprule
\textbf{Disposition} & \textbf{Count} \\
\midrule
Verified (kernel-accepted) & 179 \\
Unformalized (3 attempts exhausted) & 4 \\
Blocked (dependency not verified) & 18 \\
Excluded a priori (\texttt{thm:faithful} + n/a) & 1 \\
\midrule
Total distinct labels & 202 \\
\end{tabular}
\caption{Disposition of all 202 distinct statement labels, re-derived
from \texttt{formalization\_receipts.jsonl} (201 lines) plus the one
a-priori exclusion, which by design receives no receipt line.}
\label{tab:disposition}
\end{table}


\noindent $179 + 4 + 18 + 1 = 202$, and $202 + 17\text{ (deduplicated IRIs)} = 219$,
the full corpus IRI count. This is a disjoint partition with no
remainder category. Spelled out as an explicit chain rather than left
implicit:
\[
\begin{aligned}
219 \text{ corpus IRIs} &= 202 \text{ distinct labels} + 17 \text{ duplicate-label IRIs} \\
202 \text{ distinct labels} &= 179 \text{ verified} + 4 \text{ unformalized} + 18 \text{ blocked} + 1 \text{ excluded a priori} \\
201 \text{ receipt lines} &= 7 \text{ pilot receipts} + 194 \text{ wave-scheduled receipts}
\end{aligned}
\]
The 1 excluded statement (\texttt{thm:faithful}) receives no receipt
line by design (Section~\ref{sec:corpus-numbers}), which is why 201
receipt lines plus 1 exclusion equals 202 distinct labels, not 201.


\begin{quote}
\textbf{Accounting invariant.} \\
$\text{verified} + \text{unformalized} + \text{blocked} + \text{excluded} = \text{distinct labels}$ \\
$\text{distinct labels} + \text{duplicate-label IRIs} = \text{corpus IRIs}$ \\
$\text{receipt lines} + \text{excluded} = \text{distinct labels}$
\end{quote}
Duplicate-label IRIs are not failures: they are quotient-collapsed
before scheduling because the Lean declaration target is
label-indexed, not source-paper-IRI-indexed (Section~\ref{sec:dedup}).
A restated concept across two papers produces one Lean declaration, not
two, so counting it once is correct, not a shortfall.


Of the 202 distinct labels, 7 were already formalized and verified prior
to the wave-scheduled run reported below, as a manual pilot proving the
reject-fix-retry loop worked at all: \texttt{ax:obs},
\texttt{def:denialabstract}, \texttt{thm:rice}, \texttt{def:adm},
\texttt{def:mu}, \texttt{def:receipt}, \texttt{prop:monoid}. The
remaining 194 were processed by the automated wave scheduler described
in Chapter~\ref{ch:mathmfg}, in 8 topological waves.

\section{Lean 4 admission, not self-report}
\label{sec:lean-admission}


\begin{definition}[Lean admission]
$\mathrm{LeanAdmitted}(s)$ holds iff the \texttt{.lean} file
corresponding to statement $s$ exits $0$ under a direct invocation of
the real Lean~4 kernel:
\begin{quote}
\texttt{PATH="\$HOME/.elan/bin:\$PATH" lean tools/paper-factory/lean-pilot/<statement>.lean}
\end{quote}
\end{definition}
\begin{definition}[No self-report]
$\mathrm{AgentReportedSuccess}(s)$ -- a formalization agent's own claim
that it succeeded -- does \textbf{not} imply $\mathrm{LeanAdmitted}(s)$.
Only a direct, independent invocation of the command above, run after
and outside of the agent's own report, establishes
$\mathrm{LeanAdmitted}(s)$.
\end{definition}
Every \texttt{verified} disposition in Table~\ref{tab:disposition} was established this
way: Lean~4 is the admission court for whether a formal file
type-checks, not a semantic oracle for whether that file fully captures
an author's intended meaning (Section~\ref{sec:rice-quarantine} makes
this distinction precise). Section~\ref{sec:no-self-report} states the
same principle where the pipeline is first defined; it is restated here,
next to the actual disposition table, so the two are visually tied
rather than left to a reader's memory of an earlier chapter.

\section{The four kernel rejections}


Four statements exhausted three real \texttt{lean} invocations without
kernel acceptance. Their kind and the terminal error are retained
verbatim in the receipts file, not paraphrased into success:
\texttt{thm:mrrsep} (theorem; 4 kernel errors on the final attempt),
\texttt{prop:bottom} (proposition; \texttt{Function expected} at a
specific line), \texttt{prop:cone} (proposition; a type mismatch),
\texttt{cor:allstages} (corollary; a type mismatch on a compound boolean
term). These are reported as genuine limits of the bare-Lean-core,
bounded-retry approach (Chapter~\ref{ch:limitations}), not as pipeline
defects.

\section{Two infrastructure bugs, found and fixed during the run}


The wave-scheduler harness itself failed twice during development, in
ways worth reporting as evidence that the pipeline's own correctness was
checked rather than assumed.


\textbf{Args-as-string.} The workflow orchestration tool's top-level
data-passing parameter arrives as a raw JSON string rather than a parsed
value in its scripting sandbox. An early version of the harness attempted
to pass the full corpus (including LaTeX statement text) through this
channel; large payloads were corrupted in transit. The fix was
structural, not a patch: only a small label/kind/dependency skeleton is
passed through the orchestrator, and each per-statement formalization
agent fetches its own statement's LaTeX text directly via a SPARQL query
it runs itself.


\textbf{Wave-scheduler precompute bug.} An early version computed every
topological wave upfront from a results map that was never updated with
live outcomes, so after wave~0 genuinely ran, the scheduler evaluated
wave~1's readiness against stale (empty) results and force-blocked all
remaining statements. The fix interleaves computation and execution:
wave $n$'s readiness is computed from the live, continuously updated
results map immediately before wave $n$ executes, not from a value
computed once at the start. The corrected run was resumed from cache
(the tool's \texttt{resumeFromRunId} mechanism), so the 54 already-valid
wave-0 results were reused rather than recomputed.

\section{The receipts-consolidation gap}
\label{sec:receipts-gap}


After the wave-scheduled run completed, the consolidated
\texttt{formalization\_receipts.jsonl} was found to contain 199 lines,
not the expected 201 (7 pilot $+$ 194 processed): the merge step that
combined the manual-pilot receipts with the wave-run receipts had
silently dropped the standalone entries for \texttt{def:receipt} and
\texttt{prop:monoid}, even though both statements' \texttt{.lean} files
were verified and present on disk. This was caught by the same
independent-verification discipline stated in
Section~\ref{sec:no-self-report}: re-running the real \texttt{lean}
kernel directly against both files confirmed both type-check (the
\texttt{elan}-managed \texttt{lean} binary invoked directly on
\texttt{def\_receipt.lean} and \texttt{prop\_monoid.lean}, both exit~0),
and the
two missing receipt lines were re-added, bringing the file to 201 lines
and the counts in this chapter to their final, reconciled values.


\begin{finding}
A consolidation step (merging two receipt sources into one file) is
itself a point where evidence can be silently lost even when the
underlying artifacts remain correct on disk. The fix here was procedural
(independent kernel re-verification of every file, not just trusting the
consolidated file's line count) rather than a change to the pipeline's
formalization logic.
\end{finding}

\chapter{Civilizational-impact analysis}
\label{ch:civilizational}


Every claim in this chapter is marked with the \textsc{Speculation}
environment defined in the preamble. None of it is derived from the
219-statement result in the sense that Chapters~\ref{ch:mathmfg}--\ref{ch:empirical}
are derived: it is motivated by that result, and it is falsifiable in
principle, but it is not kernel-checked the way a Lean statement is.
Chapter~\ref{ch:conclusion}'s verification step confirms mechanically
that no claim in this chapter escapes this marking.

\section{An anchor number, not an estimate}


Before speculating, one real number, taken directly from git commit
timestamps, not from memory or narrative: the commit that added the
paper-factory tooling and compiled a validated 41-page mega-paper PDF
(\texttt{55e4cf4}, 2026-07-05 01:27:30) and the commit recording 172
newly kernel-verified Lean statements across an 8-wave scheduled run
(\texttt{8d15996}, 2026-07-05 02:26:31) are 59 minutes apart. In that
interval: 13 groups of genuine shared mathematical notation were merged
and 2 accidental symbol collisions were resolved at the source; the
whole corpus was re-validated (SHACL, dangling-citation, notation-
conflict); and a topologically-scheduled, kernel-verified formalization
run produced 172 new \texttt{.lean} artifacts, each independently
re-checked against the real Lean kernel.


\begin{speculative}
A domain expert manually selecting, drafting, and Lean-checking a single
nontrivial formalization --- reading the source statement, choosing
encodings for its dependencies, writing tactics, iterating against
kernel errors --- rarely takes less than thirty minutes even for a
routine lemma, and can take hours for a genuinely novel argument. Taking
thirty minutes as a deliberately low floor, 172 statements implies at
least 86 expert-hours of equivalent manual effort, compressed into a
59-minute wall-clock window. This is not a claim that the automated
result is equivalent in depth or reliability to 86 hours of expert
attention per statement --- Chapter~\ref{ch:limitations} states plainly
where the automated pipeline's guarantees stop --- it is a claim that
the throughput \emph{ratio}, on this one data point, is large enough
that ``how many PhDs would this replace'' is the right order-of-magnitude
question to be asking, rather than an obviously rhetorical one.
\end{speculative}

\section{What changes if verified math has a manufacturing cost curve}


\begin{speculative}
Physical manufacturing did not merely make existing artisanal goods
cheaper; it changed which goods were worth producing at all, because the
marginal cost of the tenth unit stopped resembling the marginal cost of
the first. If kernel-verified mathematical statements can be produced at
a cost curve that similarly flattens with corpus size --- because each
newly verified statement enlarges the reusable, already-proved
vocabulary available to every subsequent statement's formalization
attempt, exactly as Chapter~\ref{ch:mathmfg}'s wave scheduler does ---
then the population of mathematical claims worth formally verifying
stops being limited to results whose importance justifies a human
expert's dedicated attention, and starts including large swaths of
routine, load-bearing lemmas that are currently left as ``clearly true''
prose because no one has hours to spend proving them formally.
\end{speculative}


\begin{speculative}
For mathematical-research throughput, this suggests a shift in what a
research paper's ``proof'' section is for: not the sole evidence of
correctness, but a draft that a math-manufacturing pipeline lowers into
a kernel-checked artifact as a matter of course, the way a compiler
lowers source to machine code today. The reviewer's question moves from
``is this proof convincing'' to ``does this proof compile'' for the
mechanical portions, freeing expert attention for the genuinely novel
steps --- plausibly the four kernel-rejected statements in this thesis's
own run (Section~\ref{sec:corpus-numbers}) are exactly the kind of
step that should keep demanding a human.
\end{speculative}


\begin{speculative}
For formal-verification adoption, the barrier has historically been
that formalizing existing mathematics is itself a research-grade skill,
so verified libraries (Mathlib and its counterparts) grow only as fast
as a comparatively small community of specialists can push them forward.
A pipeline of the shape in Chapter~\ref{ch:mathmfg} --- LLM drafting,
kernel admission, automated reject-fix-retry, a compounding verified-
vocabulary library --- targets exactly that bottleneck. It does not
remove the need for the specialist community; it changes what fraction
of a specialist's day goes to formalizing routine statements versus
extending the frontier.
\end{speculative}


\begin{speculative}
For mathematics education, a manufacturing-style pipeline that can
formalize a routine claim in minutes suggests a different kind of
exercise: a student's proof attempt checked by the same kernel used
here, with the same reject-fix-retry loop turned into an interactive
tutor rather than an autonomous batch process. This is a natural
extension, not a novel proposal --- interactive theorem provers have
been used pedagogically before --- but the wave-scheduled, dependency-
aware version in this thesis suggests the tutoring loop could scale to
a whole curriculum's dependency graph rather than one exercise at a
time.
\end{speculative}

\section{Replacement cost: pricing the output, the system, and the onboarding separately}
\label{sec:replacement-cost}


Section~\ref{sec:composite-chain} separates what each layer of the
chain does; this section separates what would need to be \emph{paid
for} to reproduce it, because collapsing these into one number hides
where the real cost sits. Three things price differently: (i) the
formalization output itself (179 verified, 4 rejected, 18 blocked, 17
deduped Lean statements); (ii) the manufacturing system that produced
it (the Praxis corpus, scheduler, receipts, and \texttt{doc-pack}
pipeline of Section~\ref{sec:composite-chain}); and (iii) the
onboarding burden of a team that does not yet operate inside that
system -- learning to work \emph{in} Praxis and Lean together is a
different, larger cost than either sub-skill alone.


\begin{speculative}
The following is a back-of-envelope replacement-cost estimate, not a
quoted or audited figure: every number is a stated assumption
(team size, duration, loaded hourly rate) multiplied out explicitly, so
a reader can substitute their own assumptions and recompute rather than
take the total on faith. Loaded hourly rate assumes roughly 1.4x base
salary for benefits and overhead (a commonly cited private-industry
benefits-load ratio), applied to a base rate in the \$70--\$130/hr range
for senior technical/research staff, giving loaded rates of
\$100--\$190/hr depending on seniority and scenario.
\end{speculative}


\begin{speculative}
\begin{table}[h]
\centering
\small
\begin{tabular}{@{}p{3.4cm}p{2.8cm}p{3.4cm}p{2.4cm}@{}}
\toprule
\textbf{Scenario} & \textbf{Team x duration} & \textbf{Assumed loaded rate} & \textbf{Cost range (midpoint)} \\
\midrule
Operator + LLM, existing machinery & 1 person, 40--100 hrs & \$100--\$200/hr + \textasciitilde\$300 compute & \$4.2k--\$20.5k (\textasciitilde\$8k) \\
Small team, already inside the repo & 3--5 people, 2--4 mo & \$125--\$165/hr & \$120k--\$528k (\textasciitilde\$325k) \\
External expert team, replicating from outside & 5--8 people, 6--12 mo & \$140--\$180/hr & \$672k--\$2.76M (\textasciitilde\$1.15M) \\
Institutional build from scratch & 8--12 people, 12--24 mo & \$150--\$190/hr, +10--25\% overhead & \$2.53M--\$10.9M (\textasciitilde\$4.0M) \\
\bottomrule
\end{tabular}
\caption{Back-of-envelope replacement-cost ranges by delivery model. Midpoints computed as team~$\times$~months~$\times$~160~hrs/month~$\times$~assumed loaded rate; ranges from the stated low/high assumptions, not independently derived.}
\label{tab:replacement-cost}
\end{table}
\end{speculative}


\begin{speculative}
Dividing the midpoint of each scenario by the 179 kernel-verified
statements gives a rough per-statement replacement cost: about \$45 for
the operator-plus-LLM scenario, \$1{,}815 for a small team already
inside the repo, \$6{,}436 for an external expert team, and \$22{,}203
for an institutional rebuild. The nearly 500x spread between the
cheapest and most expensive scenario is not a claim about the intrinsic
value of a verified statement; it is a claim about how much of the
total cost is onboarding into the manufacturing system versus operating
it once built -- almost all of that spread is explained by the 6--18
person-months of onboarding time assumed for teams that do not already
operate inside the Praxis/Lean/ggen stack, not by any difference in the
formalization work itself.
\end{speculative}

\section{The load-bearing caveat}


\begin{speculative}
Every claim above assumes the throughput advantage observed on this one
219-statement corpus, run against bare Lean~4 core with no mathlib,
generalizes to larger, more mathematically diverse corpora and to
libraries with an existing large body of dependencies to search and
reuse. That assumption is untested by this thesis and is the natural
target of the future work named in Chapter~\ref{ch:limitations}.
\end{speculative}

\chapter{Related work}
\label{ch:related}


This chapter situates the empirical result of Chapter~\ref{ch:empirical}
against prior autoformalization and LLM-assisted theorem-proving work.
It does not claim priority for the general idea of using a language
model to draft formal proofs and a kernel to check them --- that idea
predates this thesis by years --- it claims a specific, narrower thing:
a dependency-aware, receipt-producing, wave-scheduled harness with
automated reject-fix-retry, applied to a from-scratch bare-Lean-core
corpus with no mathlib dependency, then used to manufacture the thesis
document itself (Section~\ref{sec:doc-mfg-correspondence}), is a
comparatively underexplored point in the design space. Any external
comparison below that has not itself been independently re-run against
this repository's own corpus is cited as related work, not treated as
evidence for this thesis's own measured claims.

\section{Large-scale verified libraries}


Lean~4 itself \cite{demoura2021lean4}, Mathlib \cite{mathlib2020}, and
comparable libraries for Coq \cite{coq2023} and Isabelle/HOL
\cite{nipkow2002isabelle} represent decades of cumulative, mostly
human-driven formalization effort, and are the existing evidence that a
compounding verified-vocabulary library is possible at scale --- the
same structural property Chapter~\ref{ch:mathmfg}'s wave scheduler
exploits at a much smaller size (a 194-statement scope, not a library
of hundreds of thousands of results). This thesis's corpus deliberately
does not depend on Mathlib; every formalization in
\texttt{tools/paper-factory/lean-pilot/} is proved from bare Lean~4 core.
That is a scope limitation (Chapter~\ref{ch:limitations}), not a claim
that avoiding Mathlib is preferable in general.

\section{LLM-assisted formal proof search}


Systems that pair a large language model with an interactive theorem
prover's kernel to search for and check proofs --- including
competition-mathematics-focused systems such as AlphaGeometry, reported
to reach medal-level performance on olympiad-style geometry problems
\cite{trinh2024alphageometry} --- share this thesis's core loop: draft,
check, revise on kernel rejection. The
distinguishing features of the pipeline in Chapter~\ref{ch:mathmfg} are
(a) the input is a pre-existing, dependency-linked corpus of a research
paper series' own definitions and theorems, rather than a benchmark of
independent competition problems, and (b) scheduling is explicitly
topological over a real \texttt{dependsOn} graph, so that a
formalization's available vocabulary strictly grows with each completed
wave, rather than each attempt starting from a fixed background library.

\section{Autoformalization of existing mathematical text}


Prior work translating existing informal mathematical text (textbooks,
papers) into formal statements and proofs \cite{szegedy2020autoformalization}
establishes that the informal-to-formal lowering step itself is a
nontrivial research problem, independent of proof search. This
thesis's corpus sidesteps part of that
problem by construction: the source statements were already extracted
into structured RDF individuals with an explicit \texttt{kind} and
\texttt{dependsOn} graph (the paper-factory pipeline, predating this
thesis), rather than starting from free-form prose. The informal-to-RDF
step is therefore prior work of this project's own, not addressed fresh
here, and is itself a limitation acknowledged in
Chapter~\ref{ch:limitations}.

\section{What is not addressed}


This thesis does not attempt cross-kernel verification (checking the
same statement in more than one proof assistant, e.g. Lean and Coq
independently, for redundancy against a single kernel's own defects),
net-new theorem generation (proposing statements not already present in
the source corpus), or benchmarking against a standard autoformalization
dataset. Each is a natural extension of the harness in
Chapter~\ref{ch:mathmfg}, not attempted here for scope reasons stated in
Chapter~\ref{ch:limitations}.

\chapter{Limitations}
\label{ch:limitations}


Stated without euphemism, in the order they matter most.

\section{Bare Lean core, no mathlib}


All 219 corpus statements' primary, currently-relied-upon
verification remains against Lean~4 core with no mathlib linked, and
\texttt{thm:faithful} remains excluded a priori for the reason given
below. What changed after this was first written: a second, additive
Lake-managed lane
(\texttt{tools/paper-factory/lean-lake/}, Section~\ref{sec:praxis-lean})
now links Mathlib via its own prebuilt cache mechanism
(\texttt{lake exe cache get}, 8{,}542 files, avoiding a from-source
build), and a bounded four-file pilot within that lane demonstrates real
reuse of Mathlib content in place of from-scratch axioms
(Section~\ref{sec:praxis-lean}). This does not retroactively migrate the
183 already-verified bare-core files, and it does not close the
capability gap that excluded \texttt{thm:faithful}: the pilot's own
cryptographic hash function (\texttt{chainH}) remains axiomatized, since
no verified BLAKE3 implementation exists in Mathlib either. The
capability boundary below is therefore softened, not eliminated: mathlib
is no longer categorically absent from this thesis's toolchain, but it is
linked in one additive lane covering four of 219 statements, not the
whole corpus.


The remaining capability boundary: a corpus leaning further into
analysis, probability, or algebraic structures with deep existing
mathlib support would still need the pilot's scope widened well past
four statements to get a fair test of the wave-scheduling and
reject-fix-retry mechanisms this thesis actually contributes, and
\texttt{thm:faithful}'s cryptographic/probabilistic content specifically
still has no pre-built Lean formalization to cite (unlike
\texttt{thm:rice}'s Section~\ref{sec:praxis-lean} pilot, which had one
sitting directly in \texttt{Mathlib.Computability.Halting}).

\section{Single kernel, no cross-check}


Every statement is checked against exactly one proof assistant's kernel.
A kernel defect (a soundness bug in Lean~4 itself) would not be caught by
this pipeline, since there is no independent second kernel (Coq,
Isabelle) checking the same statement. This is a real, if low-probability,
gap: production-grade formal-verification claims at civilizational scale
(Chapter~\ref{ch:civilizational}) would reasonably want cross-kernel
redundancy for anything load-bearing.

\section{Scope: 194 of 219, not the whole corpus}


The wave-scheduled run covers 194 of 219 statement instances (202 unique
labels minus the 7 already-verified pilot statements and 1 a-priori
exclusion, doubled back through Chapter~\ref{ch:empirical}'s accounting).
Spelled out: 219 corpus IRIs collapse to 202 distinct labels (17
duplicate-label IRIs deduplicated); of those 202, 201 received a
formalization receipt line (7 already-verified pilot statements plus 194
wave-scheduled statements) and 1 (\texttt{thm:faithful}) was excluded a
priori and therefore never received a receipt line at all -- the two
tables in Chapter~\ref{ch:empirical} agree on this chain by construction,
not by coincidence.
The remaining 18 blocked and 4 rejected statements are recorded honestly
as not-yet-verified, not silently omitted, but they are also not
verified: any claim about "the corpus" being formally checked must state
its denominator explicitly rather than leave it implicit -- 179/219
corpus IRIs ($\approx$ 82\%) or equivalently 179/202 distinct labels
($\approx$ 89\%) -- and never simply "the corpus."

\section{The receipts-consolidation gap, twice}


Section~\ref{sec:receipts-gap} already reports the first consolidation
gap: two verified statements silently dropped from a merged receipts
file, caught by independent re-verification rather than trusting the
consolidated file's own line count. This thesis document's own
authoring process surfaced a second, structurally similar risk: an
early draft of the same partition, written as prose in a conversational
exchange before this document existed, itself summed to 218 rather than
219 for an unrelated reason (a miscounted category, not the dropped-
receipt bug) -- caught and corrected before this thesis's own Chapter~4
was drafted. Two independent instances of
the same failure mode --- a consolidation or restatement step silently
dropping or miscounting verified facts --- is itself evidence that this
failure mode is not a one-off, and that any civilizational-scale version
of this pipeline (Chapter~\ref{ch:civilizational}) needs the
independent-reverification discipline built in as a structural
requirement, not an incidental good practice.

\section{The correspondence in Chapter~\ref{ch:mathmfg} is an analogy, not a proof}


Table~\ref{tab:correspondence}'s term-for-term mapping between the two
$\muop$-pipelines is a structural observation about shared shape
(admission gate, defect taxonomy, receipt), not a formal proof that the
two pipelines are instances of one algebraic construction in the sense
Chapter~\ref{ch:background}'s admission algebra makes precise for a
single pipeline. Making that correspondence itself a checked
mathematical statement --- proving $\muop$-pipeline-hood as a typeclass
or trait both implementations satisfy, say --- is future work, not
something this thesis establishes.


Stated as a concrete future-work target rather than left vague, the
typeclass this thesis does not prove would take the following shape
(pseudocode, not yet formalized in this repository):
\begin{quote}
\begin{verbatim}
structure MuPipeline where
  Raw      : Type
  Admitted : Type
  Artifact : Type
  Refusal  : Type
  admit       : Raw -> Sum Refusal Admitted
  manufacture : Admitted -> Artifact
  receipt     : Artifact -> Receipt
\end{verbatim}
\end{quote}
The corresponding proposition targets are: (i) \texttt{src/mfg.rs} and
the Lean~4 formalization harness (Chapter~\ref{ch:mathmfg}) should each
instantiate \texttt{MuPipeline}; (ii) the \texttt{doc-pack}/ggen/LaTeX
document pipeline (Section~\ref{sec:doc-mfg-correspondence}) should
instantiate the same structure at the document-artifact layer. This
thesis establishes three \emph{operational} instances and names their
common structure in prose and tables; it does not construct or
type-check the \texttt{MuPipeline} structure above or prove the three
instance obligations. That construction is the precise next theorem
target this section leaves open, not a concession that weakens the
correspondence already shown.

\section{This document's own manufacturing is new and untested at scale}


The thesis document itself (Chapters 1--8) is rendered from RDF paragraph
individuals via the \texttt{doc-pack} ggen pack, a mechanism built
specifically for this document and not previously exercised on a
document this size. The \texttt{doc:register} field
(evidence-bound/speculative) is checked mechanically in
Chapter~\ref{ch:conclusion}'s verification step for exactly the
paragraphs authored this way; it is a real, novel piece of
infrastructure, not a mature, independently-tested tool.

\chapter{Conclusion}
\label{ch:conclusion}

\section{What was actually shown}


A 219-statement mathematics corpus was routed through an automated
pipeline into a disjoint, receipted partition: 179 statements
kernel-verified against Lean~4, 4 honestly rejected after bounded retry,
18 correctly withheld on an unresolved dependency, and 17 cross-paper
restatements deduplicated before scheduling. This pipeline is a second
working instance of a manufacturing morphism already named in this
project's own prior work ($\Act = \muop(\Adm)$), alongside the existing
RDF-to-PDDL pipeline in \texttt{src/mfg.rs}. Two real infrastructure bugs
and one receipts-consolidation gap were found and fixed during the work,
each caught by independent re-verification against the real Lean kernel
or the real corpus, never by trusting a self-report.

\section{What was additionally shown, about this document itself}
\label{sec:speculation-boundary}


This document's own Chapters~1--8 are not hand-authored \texttt{.tex}:
they are rendered by \texttt{ggen sync} from RDF paragraph individuals
(\texttt{docs/thesis/math\_manufacturing/rdf/thesis.ttl}) through the
\texttt{doc-pack} ontology and template, following the same
admit/render/receipt shape as the math-manufacturing pipeline it
describes. Each paragraph in that data carries a \texttt{doc:register}
value, \texttt{evidence-bound} or \texttt{speculative}, and this is a
mechanically checkable fact, not a typographic convention. This thesis
is therefore, in a narrow and literal sense, an instance of the pattern
it argues for: its own structure is admitted data, rendered through a
manufacturing morphism, not asserted prose -- stated formally as the
third instantiation of Definition~\ref{def:mu-pipeline} in
Section~\ref{sec:doc-mfg-correspondence}, not merely asserted here.


The check itself is the following SPARQL query, run against
\texttt{thesis.ttl}:
\begin{quote}
\begin{verbatim}
SELECT ?order ?register (COUNT(?p) AS ?n) WHERE {
  ?p a doc:Paragraph ; doc:register ?register .
  { ?p doc:inChapter ?chapter }
  UNION
  { ?p doc:inSection ?sec . ?sec doc:inChapter ?chapter }
  ?chapter doc:order ?order .
} GROUP BY ?order ?register ORDER BY ?order
\end{verbatim}
\end{quote}
Its expected result, restated here as evidence rather than left to a
reader's own run: Chapters~1, 2, 3, 4, 6, 7, and 8 report
\texttt{evidence-bound} for every one of their 6, 8, 23, 13, 5, 8, and 7
paragraphs respectively (70 paragraphs total); Chapter~5 reports 3
\texttt{evidence-bound} framing paragraphs and 9 \texttt{speculative}
paragraphs, and no other chapter reports any \texttt{speculative} count
at all.
No paragraph in Chapters~1--4 or~6--8 carries \texttt{speculative}; only
Chapter~5 does, alongside its own framing paragraphs, which remain
\texttt{evidence-bound}. That is the expected result this query is
required to reproduce after any future edit to \texttt{thesis.ttl}.

\section{What was not shown}


The speculative chapter's central comparison --- that the observed
throughput on this one corpus generalizes to larger, more diverse
corpora, or that a manufacturing-style cost curve for verified
mathematics follows from one 59-minute data point --- is not shown. It
is stated as speculation, marked as such in both the LaTeX environment
and the underlying RDF fact, and left as the natural target for future
work.

\section{Adversarial review surface}
\label{sec:adversarial-review-surface}


This section does not anticipate or answer hostile arguments in prose.
It states, for each class of challenge a reviewer might raise, exactly
what evidence would have to exist for that challenge to be an admitted
refutation of one of this thesis's evidence-bound claims, rather than a
rhetorically interesting but unsubstantiated objection.


\begin{table}[h]
\centering
\small
\begin{tabular}{@{}p{4.2cm}p{7.8cm}@{}}
\toprule
\textbf{Challenge form} & \textbf{Required evidence} \\
\midrule
Count error & a failing corpus/receipt reconciliation against the accounting invariant of Section~\ref{sec:corpus-numbers} \\
False Lean success & the corresponding \texttt{.lean} file fails a direct kernel check (Section~\ref{sec:lean-admission}) \\
Bad blocker & the dependency graph shows no unresolved dependency for a statement marked blocked \\
Bad dedupe & two IRIs sharing a label are shown not to be the same target concept \\
Bad exclusion & \texttt{thm:faithful} is shown to be bare-Lean-core admissible after all \\
Bad document register & the SPARQL query of Section~\ref{sec:speculation-boundary} shows a speculative paragraph outside Chapter~5 \\
Bad ggen manufacturing claim & \texttt{regenerate.sh} fails, or its output diverges from \texttt{thesis.pdf} \\
Bad $\muop$-pipeline claim & the Definition~\ref{def:mu-pipeline} instance mapping of Table~\ref{tab:composite-chain} is shown invalid for one of its rows \\
Bad economics claim & the estimate of Section~\ref{sec:replacement-cost} is found presented as evidence rather than labeled speculation \\
\bottomrule
\end{tabular}
\caption{Admitted challenge forms and the evidence each requires.}
\end{table}
A criticism that does not instantiate one of these challenge forms may
still be rhetorically interesting, but it is not an admitted refutation
of this thesis's evidence-bound claims.

\section{Closing}


The correct adversarial question after this result is not whether gaps
exist --- Chapter~\ref{ch:limitations} names six of them plainly --- but
which specific admitted statement, blocked dependency, or
evidence-bound claim is wrong, checkable against the receipts, the
corpus, and the Lean kernel that produced it.


